% Normative source: sections 1.2, 2, 3.5, 9, 10, 13--17.
\section{논의와 한계}
\label{sec:limits}

이 절은 앞의 검증·추적·차단 결과가 무엇을 뜻하는지,
그 결과를 실제 운영에 적용하려면 무엇이 더 필요한지 구분한다.

\subsection{회계 값의 해석}

zkDPP는 원장이 수락한 값과 디지털 전이가 정해진 규칙을 따르는지
검증한다. 최초 증빙의 진실성, 실제 공정·재료·제품과 문서의 대응은
등록 승인자나 외부 검사에 달려 있다. 탄소 값도 수용한 초기 값과
운송·공정 증분의 배분 결과이지, 전체 생애주기의 측정값이 아니다.
초기 탄소를 0으로 둔 경우와 폐기물의 $a=e=0$도 실제 배출이
없었다는 뜻은 아니다. 공개 DPP 검증은 지정한 필드와 발급 기록의
일치를 확인하며 문서 전체나 규제 준수를 인증하지 않는다.

합계가 맞아도 제품별 탄소 배분이 실제 배출량을 뜻하는 것은 아니다.
식~\eqref{eq:process-carbon}은 질량 0 출력에 탄소를 배분하지 않고
정수 잔여를 양의 질량 출력에 준다. 예를 들어 탄소 합계가 1인 입력을
질량 $(0,1,1)$의 세 일반용 출력으로 나누면 탄소는 $(0,1,0)$이다.
이는 정해진 정수 회계 규칙을 따른 결과일 뿐, 두 제품의 실제
탄소발자국을 측정한 결과는 아니다. 최초 탄소 값과 공정 증분의
진실성, 제품별 배분 기준의 적절성은 여전히 외부 검증이 필요하다.
\ref{sec:model-validation}절에서는 질량 0 출력에 잔여를 주려는
이전 사례를 현행 규칙이 거부하는지 검사한다.

\subsection{감사·차단의 범위}

후속 추적은 도달한 거래의 모든 출력을 포함하는 디지털 의존관계
정책이다. 실제 오염 범위나 피해 정도를 판정하지 않으며, 시작점과
차단 승인도 신뢰된 감사자의 판단이다. 원장의 동결은 현재 미소비
Note와 미해결 Voucher의 다음 사용을 거부한다. 이미 소비한 기록을
되돌리거나 발급된 Claim을 취소하거나 물리적 유통을 막지는 못한다.
감사 중 새 파생 객체가 생기면 새 시점에서 다시 조사해야 한다.

zkDPP는 부모 참조와 출력의 미래 식별자를 한 목록에 묶어 암호화해
기록마다 바로 찾도록 한다. 다른 방법도 있다. 관련 출력과 복호화 키를 모두
모아 소비 식별자에서 출력의 종류·위치를 찾는 조회표를 만들면
부모 참조를 기록에 암호화해 두지 않아도 선행 객체를 찾을 수 있다.
단, 자료의 누락·중복을
검사해야 하며, 미소비 객체를 소유자 협력 없이 차단하려면 미래
식별자를 얻는 별도 수단도 필요하다. 직접 기록과 전체 조회표 중
어느 쪽의 비용이 낮은지는 여기서 측정하지 않았다.

\subsection{구현과 평가의 한계}
\label{sec:implementation-limits}

일반 Process 설계와 확인한 구현의 범위는 다르다.
소스로 확인한 고정 구현은 입력 3개를 일반용·폐기물 출력 하나씩으로
바꾸는 특정 정책이며, 일반용 출력의 질량은 양수여야 한다.
폐기물 출력에는 질량 0을 허용한다. 입출력 수를 바꾸는 실험용 회로는
별도 코드이며 모든 출력을 양의 질량인 일반용으로 제한한다.
따라서 이들 구현이 일반 설계의 여러 일반용 출력에 대한 질량 0
배분 규칙까지 구현한 것은 아니다~\cite{zkdpp-implementation-profile-2026}.
이 소스 확인과 앞 절의 성능 보고값을 같은 실행 결과로 간주하지 않는다.

첫째, 세 보안 성질의 암호 가정이 실제 증명계·군·해시와 인코딩에서
성립하는지, 감사키의 생성·분산·복원과 조각 전달이 올바른지 확인해야 한다.
정리 3은 공개 조각 확인 자료와 K개 미만 조각을 포함하지만,
DDH·독립 무작위 오라클·다중 정리 영지식을 전제로 한다.
참여자 부패나 능동 질의까지 다루는 결과는 아니다.

둘째, 인용한 보존 기능 시험의 대표 실행을 넘어 현행
회로·계약·클라이언트를 같은 경로에서 시험해야 한다.
아홉 거래의 정상·거부 사례, 생성 때 남긴 식별자와 실제 소비의
일치, 분기·합류·자료 누락의 추적, 전체 대상의 부분 실패·경합과
공통 시점 완료 판정을 함께 확인해야 한다.

셋째, 측정 코드·설정과 원시 결과를 연결해 독립 실행의 시간 변동을
확인해야 한다. 키 복원의 1,000회 내부 평균은 독립 반복이 아니다.
Process의 기준값 사이 21,318 gas 차이 역시 설명되지 않았다.
외부 시스템의 성능 순위를 정하려면 작업 단위와 환경을 맞춰야 한다.
